\section{Usando o Build Cache}

\begin{frame}
  \frametitle{O que é o Build Cache?}
  A cada \texttt{docker build}, o Docker verifica se pode reutilizar
  o resultado de um build anterior para cada instrução.

  \vspace{0.5em}
  \begin{itemize}
    \item Cache encontrado: Reutiliza o resultado, sem re-executar
    \item Builds mais rápidos e eficientes
  \end{itemize}
\end{frame}

\begin{frame}
  \frametitle{Quando o Cache é Invalidado?}
  \begin{itemize}
    \item Mudança no comando de uma instrução \texttt{RUN}
    \item Alteração nos arquivos copiados com \texttt{COPY} ou \texttt{ADD}: Conteúdo ou permissões
  \end{itemize}

  \vspace{0.5em}
  \begin{alertblock}{Atenção}
    Uma vez que uma camada é invalidada, \textbf{todas as camadas subsequentes}
    também são invalidadas.
  \end{alertblock}
\end{frame}

\begin{frame}
  \frametitle{O Problema do Dockerfile Ingênuo}
  \texttt{COPY . .}

  \texttt{RUN yarn install --production}

  \vspace{0.5em}
  \begin{itemize}
    \item Primeiro build: $\approx$ 20 segundos
    \item Rebuild sem alterações: $\approx$ 1 segundo (cache)
    \item Qualquer mudança no código-fonte invalida o \texttt{COPY . .}
          — forçando o \texttt{yarn install} a rodar novamente, mesmo sem
          mudanças nas dependências
  \end{itemize}
\end{frame}

\begin{frame}
  \frametitle{Dockerfile Otimizado}
  Copiar primeiro só os arquivos de dependência, depois o restante:

  \vspace{0.3em}
  \texttt{COPY package.json yarn.lock ./}

  \texttt{RUN yarn install --production}

  \texttt{COPY . .}

  \vspace{0.5em}
  \begin{itemize}
    \item Criar \texttt{.dockerignore} com \texttt{node\_modules}
          para evitar copiar dependências locais
  \end{itemize}
\end{frame}

\begin{frame}
  \frametitle{Resultado da Otimização}
  Alterando apenas o código-fonte:

  \vspace{0.5em}
  \begin{itemize}
    \item \texttt{COPY package.json} e \texttt{RUN yarn install} — permanecem em cache
    \item Rebuild cai de \textbf{16 segundos} para $\approx$ \textbf{1,2 segundo}
  \end{itemize}

  \vspace{0.5em}
  \begin{block}{Princípio Geral}
    Instruções que mudam com \textbf{menor frequência} ficam no topo —
    instruções que mudam \textbf{mais frequentemente} ficam por último.
  \end{block}
\end{frame}